\renewcommand{\thetheorem}{17.a.\arabic{theorem}}
\setcounter{theorem}{0}
\setcounter{chapter}{16}
% ---------------------------------------------------------------------------
% Provenance of the prose in this file.
%   % Extractor: <model>   the block below was transcribed from Prof. Biran's
%                          handwritten notes by that model
%   % Creator:   <model>   the block below was written by that model and has no
%                          counterpart in the notes
% A tag holds until the next tag.
% ---------------------------------------------------------------------------

\chapter{Matrices}

% Extractor: Gemini 3.1 Pro
\section{Matrix Multiplication}
\label{sec:matrix_multiplication}

% def:17.a.1
\begin{definition}[Matrix Multiplication]
\label{def:matrix_multiplication}
Let $A \in \M_{m \times n}(K)$ and $B \in \M_{n \times p}(K)$. We define a new matrix, which is called the \newterm{product} (or \newterm{multiplication}) of $A$ and $B$, denoted by $A \cdot B \in \M_{m \times p}(K)$.

Write $A = (a_{ij})_{\substack{1 \le i \le m \\ 1 \le j \le n}}$ and $B = (b_{ij})_{\substack{1 \le i \le n \\ 1 \le j \le p}}$. The matrix $A \cdot B$ has entries $(c_{ij})_{\substack{1 \le i \le m \\ 1 \le j \le p}}$, where
\begin{equation*}
    c_{ij} := a_{i1}b_{1j} + \dots + a_{in}b_{nj} = \sum_{k=1}^{n} a_{ik}b_{kj}
\end{equation*}
We often write $AB$ instead of $A \cdot B$. 
\end{definition}

\begin{importantremark}
The product $AB$ is defined only when the number of columns of $A$ equals the number of rows of $B$.
\end{importantremark}

An alternative description of $A \cdot B$ can be visualized by taking the dot product of the $i$-th row of $A$ with the $j$-th column of $B$:
\begin{equation*}
    c_{ij} = (\text{row } i \text{ of } A) \cdot (\text{column } j \text{ of } B) = a_{i1}b_{1j} + \dots + a_{ik}b_{kj} + \dots + a_{in}b_{nj}
\end{equation*}
\begin{center}
\begin{tikzpicture}[>=Stealth, thick, 
    highlight/.style={opacity=0.3},
    rowcolor/.style={fill=blue!40},
    colcolor/.style={fill=red!40},
    rescolor/.style={fill=purple!50}]

    % Matrix A (m x n) - Tall and thin
    % m = 2.5, n = 1.0
    \draw (0,0) rectangle (1.0, 2.5);
    \node at (0.5, 2.8) {$A$};
    
    % Highlight row i in A
    \fill[rowcolor] (0, 1.6) rectangle (1.0, 1.9);
    \draw[blue] (0, 1.6) rectangle (1.0, 1.9);
    \node[blue] at (0.5, 1.75) {\scriptsize row $i$};

    % Matrix B (n x p) - Short and wide
    % n = 1.0, p = 2.5
    \draw (2.2, 0) rectangle (4.7, 1.0);
    \node at (3.45, 1.5) {$B$};
    
    % Highlight col j in B
    \fill[colcolor] (3.6, 0) rectangle (3.9, 1.0);
    \draw[red] (3.6, 0) rectangle (3.9, 1.0);
    % Label col j placed above the red bar, avoiding overlap
    \node[red, anchor=south] at (3.75, 1.0) {\scriptsize col $j$};

    % Equals Sign
    \node at (5.5, 1.25) {\Large $=$};

    % Result Matrix AB (m x p) - Large square/rectangle
    % m = 2.5, p = 2.5
    \draw (6.3, 0) rectangle (8.8, 2.5);
    \node at (7.55, 2.8) {$A \cdot B$};
    
    % Highlight result cell c_ij
    \fill[rescolor] (7.7, 1.6) rectangle (8.0, 1.9);
    \draw[purple] (7.7, 1.6) rectangle (8.0, 1.9);
    \node[right=2pt, purple] at (8.8, 1.75) {\scriptsize $c_{ij}$};

    % Logical visualization arrows
    \draw[->, gray!60, dashed] (1.1, 1.75) -- (6.2, 1.75);
    \draw[->, gray!60, dashed] (3.75, -0.1) |- (6.2, 1.75);

\end{tikzpicture}
\end{center}
Another way to describe it is by considering the columns of $B$, denoted as $w_1, \dots, w_p$. The matrix $A \cdot B$ is then formed by multiplying $A$ with each column vector $w_j$:
\begin{equation*}
    A \cdot \begin{pmatrix} | & & | \\ w_1 & \dots & w_p \\ | & & | \end{pmatrix} = \begin{pmatrix} | & & | \\ A w_1 & \dots & A w_p \\ | & & | \end{pmatrix}
\end{equation*}

\begin{example}[Matrix Multiplication Example]
We compute the product of a $2 \times 3$ matrix and a $3 \times 2$ matrix:
\begin{equation*}
    \begin{pmatrix} 3 & 2 & 1 \\ 1 & 0 & 2 \end{pmatrix} \cdot \begin{pmatrix} 1 & 2 \\ 0 & 1 \\ 4 & 0 \end{pmatrix} = \begin{pmatrix} 3 \cdot 1 + 2 \cdot 0 + 1 \cdot 4 & 3 \cdot 2 + 2 \cdot 1 + 1 \cdot 0 \\ 1 \cdot 1 + 0 \cdot 0 + 2 \cdot 4 & 1 \cdot 2 + 0 \cdot 1 + 2 \cdot 0 \end{pmatrix} = \begin{pmatrix} 7 & 8 \\ 9 & 2 \end{pmatrix}
\end{equation*}
Conversely, multiplying the $3 \times 2$ matrix by the $2 \times 3$ matrix yields a $3 \times 3$ matrix:
\begin{equation*}
    \begin{pmatrix} 1 & 2 \\ 0 & 1 \\ 4 & 0 \end{pmatrix} \cdot \begin{pmatrix} 3 & 2 & 1 \\ 1 & 0 & 2 \end{pmatrix} = \begin{pmatrix} 5 & 2 & 5 \\ 1 & 0 & 2 \\ 12 & 8 & 4 \end{pmatrix}
\end{equation*}
If $A \in \M_{m \times n}(K)$ and $B \in \M_{n \times p}(K)$, then $A \cdot B$ is defined, but $B \cdot A$ is not defined, unless $p = m$.
\end{example}

\subsection{Some Special Cases}

\begin{enumerate}
    \item If $A \in \M_{m \times n}(K)$ and $v \in K_{\text{col}}^{n}$, we have previously defined the matrix-vector product $A \cdot v \in K_{\text{col}}^{m}$. However, we can also identify $K_{\text{col}}^{n} = \M_{n \times 1}(K)$. If we view $v$ as an $n \times 1$ matrix, then the matrix multiplication of $A$ and $v$ yields the exact same result as applying $A$ to the vector $v$.
    
    \item If $A \in \M_{m \times n}(K)$ and $w \in K_{\text{row}}^{m} = \M_{1 \times m}(K)$, then we can define the left-multiplication $w \cdot A \in \M_{1 \times n}(K) = K_{\text{row}}^{n}$. If $A$ is represented by its row vectors $v_1, \dots, v_m$, this operation scales and sums the rows of $A$ according to the entries of $w$.
% Extractor: Opus 5 (High)
    This gives the product a row-wise description matching the column-wise one above. If $A$ has rows $v_1, \dots, v_m$ and $B \in \M_{n \times p}(K)$ has columns $w_1, \dots, w_p$, then
    \begin{equation*}
        A \cdot B = \begin{pmatrix} | & & | \\ A w_1 & \dots & A w_p \\ | & & | \end{pmatrix} = \begin{pmatrix} \text{---} & v_1 \cdot B & \text{---} \\ & \vdots & \\ \text{---} & v_m \cdot B & \text{---} \end{pmatrix}.
    \end{equation*}
% Extractor: Gemini 3.1 Pro
    
    \item Let $v = (a_1, \dots, a_n) \in K_{\text{row}}^{n} = \M_{1 \times n}(K)$, and let $w = (b_1, \dots, b_n)\transp \in K_{\text{col}}^{n} = \M_{n \times 1}(K)$. Then the product $v \cdot w \in \M_{1 \times 1}(K) = K$ is simply the standard scalar product:
    \begin{equation*}
        v \cdot w = a_1 b_1 + \dots + a_n b_n
    \end{equation*}
\end{enumerate}

\subsection{Connections to Linear Maps}

As a conclusion from our previous discussions on linear maps, let $T: V \to W$ and $S: W \to U$ be linear maps between finite-dimensional vector spaces. Let $\mathcal{B}$, $\mathcal{C}$, and $\mathcal{E}$ be bases for $V$, $W$, and $U$, respectively. The matrix representation of the composition $S \circ T$ is given by the product of their respective representation matrices:
\[
    [S \circ T]_{\mathcal{E}}^{\mathcal{B}} = [S]_{\mathcal{E}}^{\mathcal{C}} \cdot [T]_{\mathcal{C}}^{\mathcal{B}}
\]

% Extractor: Opus 5 (High)
\begin{center}
\begin{tikzcd}[column sep=7em, row sep=3em, ampersand replacement=\&]
V \arrow[r, "T"] \arrow[rr, bend left=25, "S \circ T"] \& W \arrow[r, "S"] \& U
\end{tikzcd}
\end{center}
Here $\mathcal{B}$ is a basis of the source $V$, $\mathcal{C}$ a basis of the space $W$ in the middle, and $\mathcal{E}$ a basis of the target $U$; each matrix in the identity above carries the basis of its own source as a superscript and the basis of its own target as a subscript, and the two occurrences of $\mathcal{C}$ meet in the middle exactly where $T$ and $S$ meet.

% Extractor: Gemini 3.1 Pro

\subsection{Notation and The Identity Matrix}

Let $\mathcal{I}$ be an index set. For $i, j \in \mathcal{I}$, we define the \textbf{Kronecker-delta} as:
\begin{equation*}
    \delta_{ij} := \begin{cases} 1 & \text{if } i = j \\ 0 & \text{if } i \neq j \end{cases}
\end{equation*}
Using this, we define the \newterm{identity matrix} $I_n \in \M_{n \times n}(K)$ as $I_n = (\delta_{ij})_{\substack{1 \le i \le n \\ 1 \le j \le n}}$. Explicitly, this is:
\begin{equation*}
    I_n = \begin{pmatrix} 1 & & 0 \\ & \ddots & \\ 0 & & 1 \end{pmatrix}
\end{equation*}

\begin{remark}
\label{rem:identity_represents_id}
Let $V$ be a finite-di\-men\-sional vector space, and let $n := \dim V$. Let $\id_V: V \to V$ be the identity map, and let $\mathcal{B}$ be any basis for $V$. Then, $[\id_V]_{\mathcal{B}}^{\mathcal{B}} = I_n$. In other words, the matrix representing the identity map with respect to the basis $\mathcal{B}$ is precisely the identity matrix $I_n$.
\end{remark}

\begin{exercise}[The Identity Map is Represented by $I_n$]
\label{exc:identity_represents_id}
Prove \cref{rem:identity_represents_id}.
\end{exercise}

\subsection{Properties of Matrix Multiplication}

\begin{proposition}[Algebraic Properties of Matrix Multiplication] 
% prop:17.a.2
\label{prop:matrix_multiplication_properties}
Matrix multiplication satisfies the following algebraic properties:
\begin{enumerate}[label=\textbf{(\alph*)}]
    \item \textbf{Associativity:} For all $A \in \M_{m \times n}(K)$, $B \in \M_{n \times p}(K)$, and $C \in \M_{p \times q}(K)$, we have:
    \begin{equation*}
        (AB)C = A(BC) \in \M_{m \times q}(K)
    \end{equation*}
    \item \textbf{Distributivity:} For all $A, B \in \M_{m \times n}(K)$, $C \in \M_{n \times p}(K)$, and $C' \in \M_{q \times m}(K)$, we have:
    \begin{align*}
        (A + B)C &= AC + BC \in \M_{m \times p}(K) \\
        C'(A + B) &= C'A + C'B \in \M_{q \times n}(K)
    \end{align*}
    \item \textbf{Identity:} For all $A \in \M_{m \times n}(K)$, we have $I_m A = A$ and $A I_n = A$.
    \item \textbf{Scalar Multiplication:} For all $\alpha \in K$, $A \in \M_{m \times n}(K)$, and $B \in \M_{n \times p}(K)$, we have:
    \begin{equation*}
        \alpha(AB) = (\alpha A)B = A(\alpha B)
    \end{equation*}
\end{enumerate}
\end{proposition}

\begin{proof}
Let $T_A: K^n \to K^m$, $T_B: K^p \to K^n$, and $T_C: K^q \to K^p$ be the linear maps defined by left-multiplication with the matrices $A$, $B$, and $C$, respectively (i.e., $T_A(v) = A \cdot v$ for all $v \in K^n$, $T_B(w) = B \cdot w$ for all $w \in K^p$, and $T_C(u) = C \cdot u$ for all $u \in K^q$). For any dimension $r$, we denote the standard basis of $K^r$ by $\mathcal{E}_r$. We have already established that:
\begin{equation*}
    [T_A]_{\mathcal{E}_n}^{\mathcal{E}_m} = A, \quad [T_B]_{\mathcal{E}_p}^{\mathcal{E}_n} = B, \quad \text{and} \quad [T_C]_{\mathcal{E}_q}^{\mathcal{E}_p} = C
\end{equation*}

\begin{enumerate}[label=\textbf{(\alph*)}]
    \item We rely on the associativity of function composition, which states that $(T_A \circ T_B) \circ T_C = T_A \circ (T_B \circ T_C)$. These are both maps from $K^q$ to $K^m$. We now pass to the matrix representation of these maps with respect to the standard bases $\mathcal{E}_q$ and $\mathcal{E}_m$:
    \begin{equation*}
        [(T_A \circ T_B) \circ T_C]_{\mathcal{E}_q}^{\mathcal{E}_m} = [T_A \circ (T_B \circ T_C)]_{\mathcal{E}_q}^{\mathcal{E}_m}
    \end{equation*}
    By applying the rule for the matrix representation of composed maps, we get:
    \begin{equation*}
        [T_A \circ T_B]_{\mathcal{E}_p}^{\mathcal{E}_m} \cdot [T_C]_{\mathcal{E}_q}^{\mathcal{E}_p} = [T_A]_{\mathcal{E}_n}^{\mathcal{E}_m} \cdot [T_B \circ T_C]_{\mathcal{E}_q}^{\mathcal{E}_n}
    \end{equation*}
    Expanding further, we obtain:
    \begin{equation*}
        \left([T_A]_{\mathcal{E}_n}^{\mathcal{E}_m} \cdot [T_B]_{\mathcal{E}_p}^{\mathcal{E}_n}\right) \cdot [T_C]_{\mathcal{E}_q}^{\mathcal{E}_p} = [T_A]_{\mathcal{E}_n}^{\mathcal{E}_m} \cdot \left([T_B]_{\mathcal{E}_p}^{\mathcal{E}_n} \cdot [T_C]_{\mathcal{E}_q}^{\mathcal{E}_p}\right)
    \end{equation*}
    Substituting the matrices back in yields $(A \cdot B) \cdot C = A \cdot (B \cdot C)$.
    
    \item From the properties of linear maps, we know that $(T_A + T_B) \circ T_C = T_A \circ T_C + T_B \circ T_C$. Passing to the matrix representations, we have:
    \begin{equation*}
        [(T_A + T_B) \circ T_C]_{\mathcal{E}_q}^{\mathcal{E}_m} = [T_A \circ T_C + T_B \circ T_C]_{\mathcal{E}_q}^{\mathcal{E}_m}
    \end{equation*}
    This implies that:
    \begin{equation*}
        [T_A + T_B]_{\mathcal{E}_n}^{\mathcal{E}_m} \cdot [T_C]_{\mathcal{E}_q}^{\mathcal{E}_n} = [T_A \circ T_C]_{\mathcal{E}_q}^{\mathcal{E}_m} + [T_B \circ T_C]_{\mathcal{E}_q}^{\mathcal{E}_m}
    \end{equation*}
    Which directly gives us $(A + B) \cdot C = A \cdot C + B \cdot C$. The other distributive identity in \textbf{(b)} is proven similarly.
    
    \item The identity matrix $I_m$ represents the identity map: $I_m = [\id_{K^m}]_{\mathcal{E}_m}^{\mathcal{E}_m}$. Since composing with the identity map does not change a function, we have $\id_{K^m} \circ T_A = T_A$. Therefore:
    \begin{equation*}
        [\id_{K^m} \circ T_A]_{\mathcal{E}_n}^{\mathcal{E}_m} = [T_A]_{\mathcal{E}_n}^{\mathcal{E}_m} = A
    \end{equation*}
    This implies that $I_m \cdot A = A$. The proof for $A I_n = A$ follows the same logic.
    
    \item Let $T_A: K^n \to K^m$ and $T_B: K^p \to K^n$ be the linear maps associated with $A$ and $B$, respectively. We want to show that $\alpha(AB) = (\alpha A)B = A(\alpha B)$. This is equivalent to showing that the corresponding linear maps are equal.
    
    First, consider the map $Q: K^n \to K^n$ defined by $Q(v) := \alpha v$. This map is linear. If $\mathcal{E}_n$ is the standard basis for $K^n$, then for any $e_j \in \mathcal{E}_n$, $Q(e_j) = \alpha e_j$. Thus, the matrix representation of $Q$ with respect to $\mathcal{E}_n$ is $[Q]_{\mathcal{E}_n}^{\mathcal{E}_n} = \alpha I_n$.
    
    Now, let's look at the compositions:
    \begin{itemize}
        \item The linear map corresponding to $(\alpha A)B$ is $T_{\alpha A} \circ T_B$. We know that $T_{\alpha A}(v) = (\alpha A)v = \alpha(Av) = \alpha T_A(v)$, so $T_{\alpha A} = \alpha T_A$. Thus, $T_{(\alpha A)B} = (\alpha T_A) \circ T_B$.
        \item The linear map corresponding to $A(\alpha B)$ is $T_A \circ T_{\alpha B}$. Similarly, $T_{\alpha B} = \alpha T_B$. Thus, $T_{A(\alpha B)} = T_A \circ (\alpha T_B)$.
    \end{itemize}
    
    We need to show that $(\alpha T_A) \circ T_B = T_A \circ (\alpha T_B)$. For any $v \in K^p$:
    \begin{align*}
        ((\alpha T_A) \circ T_B)(v) &= (\alpha T_A)(T_B(v)) \\
        &= \alpha (T_A(T_B(v))) \quad (\text{by definition of scalar multiplication for maps})
    \end{align*}
    And for the other side:
    \begin{align*}
        (T_A \circ (\alpha T_B))(v) &= T_A((\alpha T_B)(v)) \\
        &= T_A(\alpha (T_B(v))) \quad (\text{by definition of scalar multiplication for maps}) \\
        &= \alpha (T_A(T_B(v))) \quad (\text{by linearity of } T_A)
    \end{align*}
    Since both compositions yield the same result, the corresponding matrices are equal: $(\alpha A)B = A(\alpha B)$.
    
    Finally, to show $\alpha(AB) = (\alpha A)B$:
    The matrix $\alpha(AB)$ corresponds to the linear map $\alpha(T_A \circ T_B)$.
    The matrix $(\alpha A)B$ corresponds to the linear map $(\alpha T_A) \circ T_B$.
    We have already shown that $(\alpha T_A) \circ T_B = \alpha (T_A \circ T_B)$ (from the first part of the derivation above).
    Therefore, $\alpha(AB) = (\alpha A)B$.
\end{enumerate}
\end{proof}

% Extractor: Opus 5 (High)
The notes prepare statement \textbf{(d)} differently, through the following observation about the map that scales every vector by a fixed factor.

\begin{claim}[The Scaling Map is Represented by $\alpha I_n$]
\label{claim:scaling_map_matrix}
Let $V$ be a finite-di\-men\-sional vector space over $K$, let $\alpha \in K$, and let $\mathcal{B}$ be a basis for $V$. Consider the map $Q: V \to V$ defined by $Q(v) := \alpha v$. Then $Q$ is linear, and moreover
\[
[Q]_{\mathcal{B}}^{\mathcal{B}} = \alpha \cdot I_n = \begin{pmatrix} \alpha & & 0 \\ & \ddots & \\ 0 & & \alpha \end{pmatrix}, \qquad \text{where } n = \dim V.
\]
\end{claim}

\begin{exercise}[Preparation for Statement (d)]
\label{exc:scaling_map_and_statement_d}
Prove \cref{claim:scaling_map_matrix}, and use it to prove statement \textbf{(d)} of \cref{prop:matrix_multiplication_properties}.
\end{exercise}

\begin{remark}
One can prove \cref{prop:matrix_multiplication_properties} also by direct calculation with the entries. It is worth trying once, say for $(AB)C = A(BC)$, to see that it is not very pleasant --- and to appreciate how much work the passage to linear maps saves.
\end{remark}

% Extractor: Gemini 3.1 Pro

\begin{remark}
In the space $\M_{n \times n}(K)$, we can add and subtract matrices, and we can also multiply matrices. There is additionally a multiplicative unit $I_n$. Because these operations satisfy associativity, distributivity, and other relevant axioms, it follows that $\M_{n \times n}(K)$ forms a (non-com\-mu\-ta\-tive) ring with a unity.
\end{remark}

\begin{importantremark}
Matrix multiplication is, in general, \textbf{not commutative}. There exist matrices $A$ and $B$ such that $A \cdot B \neq B \cdot A$. For example:
\begin{align*}
    \begin{pmatrix} 1 & 0 \\ 0 & 0 \end{pmatrix} \cdot \begin{pmatrix} 0 & 0 \\ 1 & 0 \end{pmatrix} &= \begin{pmatrix} 0 & 0 \\ 0 & 0 \end{pmatrix} \\
    \begin{pmatrix} 0 & 0 \\ 1 & 0 \end{pmatrix} \cdot \begin{pmatrix} 1 & 0 \\ 0 & 0 \end{pmatrix} &= \begin{pmatrix} 0 & 0 \\ 1 & 0 \end{pmatrix}
\end{align*}
\end{importantremark}

\subsection{Invertible Matrices}

\begin{definition}[Invertibility]
An $n \times n$ matrix $A \in \M_{n \times n}(K)$ is called \newterm{invertible} if there exists a matrix $B \in \M_{n \times n}(K)$ such that $A \cdot B = B \cdot A = I_n$.
\end{definition}

\begin{remark}
If $A$ is invertible, then there exists \textbf{only one} matrix $B$ such that $A \cdot B = B \cdot A = I_n$.
\end{remark}

\begin{proof}
Assume there are two matrices $B$ and $C$ such that $AB = BA = I_n$ and $AC = CA = I_n$. Then, we can calculate:
\begin{equation*}
    B = B \cdot I_n = B \cdot (A \cdot C) = (B \cdot A) \cdot C = I_n \cdot C = C
\end{equation*}
This proves the uniqueness.
\end{proof}

\begin{definition}[Inverse Matrix and the General Linear Group]
If $A$ is invertible, the unique matrix $B$ such that $AB = BA = I_n$ is called the \newterm{inverse} of $A$ and is denoted by $A^{-1}$. We have:
\begin{equation*}
    A \cdot A^{-1} = A^{-1} \cdot A = I_n
\end{equation*}
The set of all invertible $n \times n$ matrices over $K$ is denoted by:
\begin{equation*}
    \GL_n(K) = \GL(n, K) := \{A \in \M_{n \times n}(K) \mid A \text{ is invertible}\}
\end{equation*}
\end{definition}

\begin{definition}[Group]
Let $G$ be a set and $\cdot : G \times G \to G$, $(a,b) \mapsto a \cdot b \in G$, be a binary operation. We say that $(G, \cdot)$ is a \newterm{group} if the following axioms hold:
\begin{enumerate}
    \item \textbf{Associativity:} For all $g_1, g_2, g_3 \in G$, we have $g_1 \cdot (g_2 \cdot g_3) = (g_1 \cdot g_2) \cdot g_3$.
    \item \textbf{Identity Element:} There exists an element $e \in G$ such that $e \cdot g = g \cdot e = g$ for all $g \in G$. The element $e$ is called the \newterm{unit} (or \newterm{identity}) of $G$.
    \item \textbf{Inverse Element:} For every $g \in G$, there exists an element $h \in G$ such that $g \cdot h = h \cdot g = e$. The element $h$ is called the \newterm{inverse} of $g$.
\end{enumerate}
Furthermore, if $G$ satisfies $g_1 \cdot g_2 = g_2 \cdot g_1$ for all $g_1, g_2 \in G$, we say that $G$ is a \newterm{commutative group} (or \newterm{Abelian group}).
\end{definition}

\begin{proposition}[Group Properties]
% prop:17.a.6
\label{prop:group_basic_properties}
Let $G$ be a group. Then:
\begin{enumerate}
    \item The unit $e \in G$ is unique. (The unit $e$ is sometimes denoted as $1 \in G$.)
    \item For every $g \in G$, its inverse is uniquely defined. It is denoted by $g^{-1}$ (so that $g \cdot g^{-1} = g^{-1} \cdot g = e$).
    \item For every $g \in G$, we have $(g^{-1})^{-1} = g$.
    \item For all $g_1, g_2 \in G$, we have $(g_1 \cdot g_2)^{-1} = g_2^{-1} \cdot g_1^{-1}$.
\end{enumerate}
\end{proposition}

\begin{proposition}[The General Linear Group]
% prop:17.a.7
\label{prop:gl_group_structure}
The set $\GL_n(K)$, endowed with matrix multiplication $(A, B) \mapsto A \cdot B$, is a group. The unity is the identity matrix $I_n$. The inverse of a matrix $A \in \GL_n(K)$ is $A^{-1}$. Moreover, for all $A, B \in \GL_n(K)$, we have:
\begin{equation*}
    (A^{-1})^{-1} = A \quad \text{and} \quad (AB)^{-1} = B^{-1}A^{-1}
\end{equation*}
\end{proposition}

\begin{proof}
Let $A, B \in \GL_n(K)$. To establish that $\GL_n(K)$ forms a group, we first need to show closure, namely that $AB \in \GL_n(K)$ as well. Indeed, we can verify this by checking the inverse property:
\begin{equation*}
    (B^{-1} \cdot A^{-1})(AB) = B^{-1}A^{-1}AB = B^{-1}I_nB = B^{-1}B = I_n
\end{equation*}
and, multiplying from the right,
\begin{equation*}
    (AB)(B^{-1}A^{-1}) = ABB^{-1}A^{-1} = AI_nA^{-1} = AA^{-1} = I_n
\end{equation*}
This proves that $AB$ is invertible with inverse $B^{-1}A^{-1}$, and therefore, $AB \in \GL_n(K)$. The other statements in the proposition follow immediately from what we have proved earlier regarding the uniqueness of inverses.
\end{proof}

\begin{corollary}[Unique Solution via Matrix Inverse]
% cor:17.a.8
\label{cor:unique_solution_via_inverse}
Let $A \in \GL_n(K)$. Then, for all $b \in K^n$, the linear system of equations $A \cdot x = b$ has a unique solution, namely $x = A^{-1} \cdot b$.
\end{corollary}

\begin{proof}
If $Ax = b$, then multiplying both sides from the left by $A^{-1}$ yields:
\begin{equation*}
    x = I_n x = (A^{-1} A) x = A^{-1}(Ax) = A^{-1} b.
\end{equation*}
Conversely, if we substitute $x = A^{-1}b$ into the equation, we obtain:
\begin{equation*}
    A(A^{-1}b) = (AA^{-1})b = I_n b = b
\end{equation*}
This confirms both the existence and uniqueness of the solution.
\end{proof}

\begin{proposition}[Equivalence of Invertibility and Isomorphism]
% prop:17.a.9
\label{prop:invertibility_isomorphism_equivalence}
Let $A \in \M_{n \times n}(K)$. Then, $A \in \GL_n(K)$ (i.e., $A$ is invertible) if and only if the associated linear map $T_A: K^n \to K^n$ is an isomorphism. Moreover, in this case, $(T_A)^{-1} = T_{A^{-1}}$.
\end{proposition}

\begin{proof}
First, assume that $A$ is invertible. We claim that $T_{A^{-1}} \circ T_A = \id_{K^n}$ and $T_A \circ T_{A^{-1}} = \id_{K^n}$, which will show that $T_A$ is an isomorphism. Indeed, for all $v \in K^n$:
\begin{equation*}
    T_{A^{-1}} \circ T_A(v) = A^{-1} \cdot (T_A(v)) = A^{-1} \cdot A \cdot v = I_n \cdot v = v
\end{equation*}
Similarly, one shows that for all $v \in K^n$:
\begin{equation*}
    T_A \circ T_{A^{-1}}(v) = I_n \cdot v = v
\end{equation*}
This proves that $T_A$ is an isomorphism.
 
Conversely, assume that $T_A: K^n \to K^n$ is an isomorphism. Let $S := (T_A)^{-1}$. By \cref{lem:all_linear_maps_are_matrices}, we know that any linear map from $K^n$ to $K^n$ is given by left-multiplication with a matrix, so $S = T_B$ for some $B \in \M_{n \times n}(K)$. Now, for all $w \in K^n$:
\begin{equation*}
    w = S \circ T_A(w) = T_B(A \cdot w) = B \cdot A \cdot w = (BA) \cdot w
\end{equation*}
We apply this equality for the standard basis vectors $w = e_1, w = e_2, \dots, w = e_n$ and obtain that $BA = I_n$. Similarly, using $T_A \circ S = \id_{K^n}$, one shows that $AB = I_n$ too. Thus, $A$ is invertible.
\end{proof}

\begin{exercise}[Properties of the Transpose Map]
\label{exc:transpose_properties_multiplication}
\begin{enumerate}[label=\textbf{(\alph*)}]
    \item Let $A \in \M_{m \times n}(K)$ and $B \in \M_{n \times p}(K)$. Prove that $(AB)\transp = B\transp \cdot A\transp$. \newline (\textit{Hint: Proving this by direct calculation is not so bad.})
    \item Let $A \in \GL_n(K)$. Prove that $A\transp \in \GL_n(K)$ and that $(A\transp)^{-1} = (A^{-1})\transp$.
\end{enumerate}
\end{exercise}

\begin{definition}[Triangular and Diagonal Matrices]
% def:17.a.10
\label{def:triangular_diagonal_matrices}
A matrix $A \in \M_{n \times n}(K)$, with entries $A = (a_{ij})_{\substack{1 \le i \le n \\ 1 \le j \le n}}$, is called:
\begin{itemize}
    \item \newterm{Upper triangular} if $a_{ij} = 0$ for all $i > j$.
    \item \newterm{Lower triangular} if $a_{ij} = 0$ for all $i < j$.
    \item \newterm{Diagonal} if $a_{ij} = 0$ for all $i \neq j$. A diagonal matrix has the form:
    \begin{equation*}
        \begin{pmatrix} * & & 0 \\ & \ddots & \\ 0 & & * \end{pmatrix}
    \end{equation*}
\end{itemize}
\end{definition}

\begin{lemma}[Closure of Triangular Matrices]
% lemma:17.a.11
\label{lemma:triangular_multiplication_closure}
Let $A$ and $B$ be two diagonal matrices (respectively, both upper triangular, or both lower triangular). Then, their product $A \cdot B$ is of the same type.
\end{lemma}

% Extractor: Opus 5 (High)
The notes leave the proof as an exercise; it is carried out among the solutions
at the end of this section, on \cpageref{sec:solutions_matrix_multiplication}.

% Creator: Opus 5 (High)
\subsection{Solutions to the Exercises}
\label{sec:solutions_matrix_multiplication}

\begin{exercisesolution}[Proof of \cref{rem:identity_represents_id}]
Let $\mathcal{B} = (v_1, \dots, v_n)$. The $j$-th column of $[\id_V]_{\mathcal{B}}^{\mathcal{B}}$ is the coordinate vector of $\id_V(v_j) = v_j$ with respect to $\mathcal{B}$. Writing $v_j = 0 \cdot v_1 + \dots + 1 \cdot v_j + \dots + 0 \cdot v_n$ shows that this coordinate vector is $e_j$, whose $i$-th entry is $\delta_{ij}$. So the $(i,j)$-entry of $[\id_V]_{\mathcal{B}}^{\mathcal{B}}$ is $\delta_{ij}$, which is exactly $I_n$.
\end{exercisesolution}

\begin{exercisesolution}[Proof of \cref{claim:scaling_map_matrix} and of \cref{prop:matrix_multiplication_properties} \textbf{(d)}]
The map $Q(v) := \alpha v$ is linear, since $Q(a v + b v') = \alpha(a v + b v') = a \alpha v + b \alpha v' = a Q(v) + b Q(v')$. For its matrix, the $j$-th column of $[Q]_{\mathcal{B}}^{\mathcal{B}}$ is $[Q(v_j)]_{\mathcal{B}} = [\alpha v_j]_{\mathcal{B}} = \alpha e_j$, and a matrix whose $j$-th column is $\alpha e_j$ for every $j$ is precisely $\alpha I_n$.

Now let $\alpha \in K$, $A \in \M_{m \times n}(K)$ and $B \in \M_{n \times p}(K)$. Applying the claim in $K^m$ with respect to the standard basis, multiplication by $\alpha$ on $K^m$ is represented by $\alpha I_m$, and therefore $\alpha C = (\alpha I_m) \cdot C$ for every $C \in \M_{m \times q}(K)$; likewise $\alpha A = (\alpha I_m) A$. Using associativity, \cref{prop:matrix_multiplication_properties} \textbf{(a)}, twice:
\[
\alpha (AB) = (\alpha I_m)(AB) = \big((\alpha I_m) A\big) B = (\alpha A) B.
\]
For the remaining identity, note that $(\alpha I_m) A = A (\alpha I_n)$, since both sides have $(i,j)$-entry $\alpha a_{ij}$. Hence
\[
(\alpha A) B = \big(A (\alpha I_n)\big) B = A \big((\alpha I_n) B\big) = A (\alpha B),
\]
again by associativity. Together these give $\alpha(AB) = (\alpha A)B = A(\alpha B)$.
\end{exercisesolution}

\begin{exercisesolution}[Proof of \cref{exc:transpose_properties_multiplication}]
\begin{enumerate}[label=\textbf{(\alph*)}]
    \item Write $A = (a_{ij}) \in \M_{m \times n}(K)$ and $B = (b_{jk}) \in \M_{n \times p}(K)$, and put $C := AB$, so that $c_{ik} = \sum_{\ell=1}^{n} a_{i\ell} b_{\ell k}$. Both $(AB)\transp$ and $B\transp A\transp$ lie in $\M_{p \times m}(K)$, so it suffices to compare entries. The $(k,i)$-entry of $(AB)\transp$ is
    \[
    c_{ik} = \sum_{\ell=1}^{n} a_{i\ell} b_{\ell k}.
    \]
    On the other hand, the $(k,i)$-entry of $B\transp A\transp$ is the $k$-th row of $B\transp$ times the $i$-th column of $A\transp$, that is,
    \[
    \sum_{\ell=1}^{n} (B\transp)_{k\ell} (A\transp)_{\ell i} = \sum_{\ell=1}^{n} b_{\ell k} a_{i \ell},
    \]
    which is the same sum, since the entries are scalars and $K$ is commutative. Hence $(AB)\transp = B\transp A\transp$.
    \item Let $A \in \GL_n(K)$. Using \textbf{(a)} and $I_n\transp = I_n$,
    \[
    A\transp \cdot (A^{-1})\transp = (A^{-1} A)\transp = I_n\transp = I_n, \qquad (A^{-1})\transp \cdot A\transp = (A A^{-1})\transp = I_n\transp = I_n.
    \]
    So $A\transp$ has a two-sided inverse, which means $A\transp \in \GL_n(K)$ and $(A\transp)^{-1} = (A^{-1})\transp$. Note that no determinants are needed: part \textbf{(a)} alone produces the inverse explicitly.
\end{enumerate}
\end{exercisesolution}

\begin{exercisesolution}[Proof of \cref{lemma:triangular_multiplication_closure}]
We give the upper triangular case; the other two follow from it.

Let $A = (a_{ij})$ and $B = (b_{ij})$ be upper triangular, so that $a_{ik} = 0$ whenever $i > k$, and $b_{kj} = 0$ whenever $k > j$. Write $C := AB = (c_{ij})$, so that $c_{ij} = \sum_{k=1}^{n} a_{ik} b_{kj}$.

Fix a pair of indices with $i > j$, and consider a single term $a_{ik} b_{kj}$ of that sum. Exactly one of the following two cases occurs:
\begin{itemize}
    \item $k < i$. Then $a_{ik} = 0$, because $A$ is upper triangular.
    \item $k \ge i$. Then $k \ge i > j$, so $k > j$, and hence $b_{kj} = 0$, because $B$ is upper triangular.
\end{itemize}
In either case the term vanishes, so $c_{ij} = 0$ for all $i > j$, which says precisely that $C$ is upper triangular.

For the lower triangular case, note that $A$ is lower triangular if and only if $A\transp$ is upper triangular. So if $A$ and $B$ are lower triangular, then $A\transp$ and $B\transp$ are upper triangular, and by \cref{exc:transpose_properties_multiplication} \textbf{(a)} together with the case already proved,
\[
(AB)\transp = B\transp A\transp \quad \text{is upper triangular,}
\]
which makes $AB$ lower triangular. Finally, a matrix is diagonal exactly when it is both upper and lower triangular, so the diagonal case follows from the two cases just proved.
\end{exercisesolution}

\begin{ainote}
Every numbered statement of the notes for this section is present in the
transcript, and 17.a.1--17.a.11 already matched the handwritten numbering. What
was missing sits around them.

Restored: the preparation the notes give for statement \textbf{(d)} of
\cref{prop:matrix_multiplication_properties}, namely
\cref{claim:scaling_map_matrix}, that scaling by $\alpha$ is represented by
$\alpha I_n$, which is how the notes ask for that statement to be proved; the
remark that the proposition can also be proved by direct calculation, with Prof.
Biran's warning that it is \qt{not very pleasant}; and the row-wise description
of a product in the second special case, which the notes give alongside the
column-wise one.

Removed: a second, verbatim copy of the remark that $\M_{n \times n}(K)$ is a
non-com\-mu\-ta\-tive ring. It appeared twice in the transcript, a few lines
apart, and once in the notes.

Two things were wrong rather than missing. The diagram under \qt{Connections to
Linear Maps} placed $\Sp(\mathcal{B})$ below $V$ joined by a rotated
\qt{$\in$}, asserting $\Sp(\mathcal{B}) \in V$, which is not a relation between
a span and the space containing it; the notes have no such diagram, and it is
replaced by the composition diagram with the bookkeeping said in words. And the
proof of \cref{lemma:triangular_multiplication_closure} opened with the word
\qt{Exercise} and then argued, in the middle of the paragraph, that \qt{for
$a_{ik}$ to be non-zero, $k$ must be less than or equal to $i$} --- the
condition runs the other way. The conclusion was right and the argument was
not.

The notes mark four things \qt{Exc.}: the identity matrix remark, the
preparation claim and statement \textbf{(d)}, the two transpose identities, and
\cref{lemma:triangular_multiplication_closure}. All four are now stated as
exercises and solved on \cpageref{sec:solutions_matrix_multiplication}. The
transcript had instead embedded \qt{proof sketches} inside the exercise
environment, and one of them derived $A\transp \in \GL_n(K)$ from
$\det(A\transp) = \det(A)$ --- determinants are three chapters away, and the
transpose identity of part \textbf{(a)} produces the inverse directly.

Finally, terms being defined for the first time --- the product, the
Kronecker-delta, the identity matrix, invertible, the inverse, upper and lower
triangular, diagonal --- were marked with \verb|\qt| rather than
\verb|\newterm|, which the house style reserves for exactly this.
\end{ainote}
